% GENERATED FILE. Edit canonical lessons, then run npm run generate:books.
% canonical-source-hash: fnv1a64:1bc49bb7d54d3046
% canonical-lessons: RU-C125-free, RU-C125-open, RU-C125-closed, RU-C125-quiet, RU-C125-noisy, RU-R125-first-pass-animals-the-body-and-the-table, RU-R125-second-pass-the-house-and-describing-words

\chapter{Free, Open, Closed, Quiet, Noisy}
\label{ch:ru-free-open-closed-quiet-noisy}
\hlchaptermodality{125}

\begin{chapteropening}
\textbf{By the end of this chapter:} \emph{I can say free, open, closed, quiet and noisy.}

Complete the last lesson of chapter 125: i can say free, open, closed, quiet and noisy.
\end{chapteropening}

\section[svobódnyy]{\ru{свободный} (svobódnyy) — free}

\label{lesson:RU-C125-free}



\begin{quote}
Before the new one: say the Russian for orange, then the Russian for busy.
\end{quote}

\subsection*{You'll want to know: \ru{свободный}}
\textbf{\ru{свободный}} (\emph{svobódnyy}) — \textquotedblleft{}free\textquotedblright{}.

\begin{grammarlens}[title={What you've built}]
One more word for this chapter.
\end{grammarlens}

\subsection*{Your turn}
\begin{itemize}
\raggedright
  \item \emph{Say it:} \emph{svobódnyy}
  \item \emph{Say it:} \emph{svobódnyy}, once more
  \item \emph{Say it:} say \emph{svobódnyy}
  \item \emph{Recall:} say the Russian for orange, then the Russian for busy, then say \emph{svobódnyy} again
\end{itemize}

\begin{tcolorbox}[breakable,colback=teal!4,colframe=teal!35!black,title={Before you move on}]
What does \ru{свободный} mean? (\textquotedblleft{}Free\textquotedblright{}.) Say it once more.
\end{tcolorbox}
\section[otkrýtyy]{\ru{открытый} (otkrýtyy) — open}

\label{lesson:RU-C125-open}



\begin{quote}
Before the new one: say the Russian for busy, then the Russian for free.
\end{quote}

\subsection*{You'll want to know: \ru{открытый}}
\textbf{\ru{открытый}} (\emph{otkrýtyy}) — \textquotedblleft{}open\textquotedblright{}.

\begin{grammarlens}[title={What you've built}]
One more word for this chapter.
\end{grammarlens}

\subsection*{Your turn}
\begin{itemize}
\raggedright
  \item \emph{Say it:} \emph{otkrýtyy}
  \item \emph{Say it:} \emph{otkrýtyy}, once more
  \item \emph{Say it:} say \emph{otkrýtyy}
  \item \emph{Recall:} say the Russian for busy, then the Russian for free, then say \emph{otkrýtyy} again
\end{itemize}

\begin{tcolorbox}[breakable,colback=teal!4,colframe=teal!35!black,title={Before you move on}]
What does \ru{открытый} mean? (\textquotedblleft{}Open\textquotedblright{}.) Say it once more.
\end{tcolorbox}
\section[zakrýtyy]{\ru{закрытый} (zakrýtyy) — closed}

\label{lesson:RU-C125-closed}



\begin{quote}
Before the new one: say the Russian for free, then the Russian for open.
\end{quote}

\subsection*{You'll want to know: \ru{закрытый}}
\textbf{\ru{закрытый}} (\emph{zakrýtyy}) — \textquotedblleft{}closed\textquotedblright{}.

\begin{grammarlens}[title={What you've built}]
One more word for this chapter.
\end{grammarlens}

\subsection*{Your turn}
\begin{itemize}
\raggedright
  \item \emph{Say it:} \emph{zakrýtyy}
  \item \emph{Say it:} \emph{zakrýtyy}, once more
  \item \emph{Say it:} say \emph{zakrýtyy}
  \item \emph{Recall:} say the Russian for free, then the Russian for open, then say \emph{zakrýtyy} again
\end{itemize}

\begin{tcolorbox}[breakable,colback=teal!4,colframe=teal!35!black,title={Before you move on}]
What does \ru{закрытый} mean? (\textquotedblleft{}Closed\textquotedblright{}.) Say it once more.
\end{tcolorbox}
\section[tíkhiy]{\ru{тихий} (tíkhiy) — quiet}

\label{lesson:RU-C125-quiet}



\begin{quote}
Before the new one: say the Russian for open, then the Russian for closed.
\end{quote}

\subsection*{You'll want to know: \ru{тихий}}
\textbf{\ru{тихий}} (\emph{tíkhiy}) — \textquotedblleft{}quiet\textquotedblright{}.

\begin{grammarlens}[title={What you've built}]
One more word for this chapter.
\end{grammarlens}

\subsection*{Your turn}
\begin{itemize}
\raggedright
  \item \emph{Say it:} \emph{tíkhiy}
  \item \emph{Say it:} \emph{tíkhiy}, once more
  \item \emph{Say it:} say \emph{tíkhiy}
  \item \emph{Recall:} say the Russian for open, then the Russian for closed, then say \emph{tíkhiy} again
\end{itemize}

\begin{tcolorbox}[breakable,colback=teal!4,colframe=teal!35!black,title={Before you move on}]
What does \ru{тихий} mean? (\textquotedblleft{}Quiet\textquotedblright{}.) Say it once more.
\end{tcolorbox}
\section[shúmnyy]{\ru{шумный} (shúmnyy) — noisy}

\label{lesson:RU-C125-noisy}



\begin{quote}
Before the new one: say the Russian for closed, then the Russian for quiet.
\end{quote}

\subsection*{You'll want to know: \ru{шумный}}
\textbf{\ru{шумный}} (\emph{shúmnyy}) — \textquotedblleft{}noisy\textquotedblright{}.

\begin{grammarlens}[title={What you've built}]
One more word for this chapter.
\end{grammarlens}

\subsection*{Your turn}
\begin{itemize}
\raggedright
  \item \emph{Say it:} \emph{shúmnyy}
  \item \emph{Say it:} \emph{shúmnyy}, once more
  \item \emph{Say it:} say \emph{shúmnyy}
  \item \emph{Recall:} say the Russian for closed, then the Russian for quiet, then say \emph{shúmnyy} again
\end{itemize}

\begin{tcolorbox}[breakable,colback=teal!4,colframe=teal!35!black,title={Before you move on}]
What does \ru{шумный} mean? (\textquotedblleft{}Noisy\textquotedblright{}.) Say it once more.
\end{tcolorbox}
\section[(dialogue)]{Animals, the body and the table — a first pass}

\label{lesson:RU-R125-first-pass-animals-the-body-and-the-table}



\begin{quote}
Say the words for quiet and noisy.
\end{quote}

\subsection*{Your turn}
Say these aloud:

\begin{itemize}
\raggedright
  \item a desert, a volcano, a cave, a waterfall, a donkey
  \item a rabbit, a tiger, a butterfly, a frog, a tortoise
  \item a squirrel, an owl, a whale, a dolphin, an elbow
  \item an ankle, a wrist, a cheek, lips, the chest
  \item a lion, a monkey, an elephant, a bee, an ant
\end{itemize}

\begin{tcolorbox}[breakable,colback=teal!4,colframe=teal!35!black,title={Before you move on}]
A desert? (\textbf{\ru{пустыня}}.) A volcano? (\textbf{\ru{вулкан}}.) A cave? (\textbf{\ru{пещера}}.) A waterfall? (\textbf{\ru{водопад}}.) A donkey? (\textbf{\ru{осёл}}.) A rabbit? (\textbf{\ru{кролик}}.) A tiger? (\textbf{\ru{тигр}}.) A butterfly? (\textbf{\ru{бабочка}}.) A frog? (\textbf{\ru{лягушка}}.) A tortoise? (\textbf{\ru{черепаха}}.) A squirrel? (\textbf{\ru{белка}}.) An owl? (\textbf{\ru{сова}}.) A whale? (\textbf{\ru{кит}}.) A dolphin? (\textbf{\ru{дельфин}}.) An elbow? (\textbf{\ru{локоть}}.) An ankle? (\textbf{\ru{лодыжка}}.) A wrist? (\textbf{\ru{запястье}}.) A cheek? (\textbf{\ru{щека}}.) Lips? (\textbf{\ru{губы}}.) The chest? (\textbf{\ru{грудь}}.) A lion? (\textbf{\ru{лев}}.) A monkey? (\textbf{\ru{обезьяна}}.) An elephant? (\textbf{\ru{слон}}.) A bee? (\textbf{\ru{пчела}}.) An ant? (\textbf{\ru{муравей}}.)
\end{tcolorbox}
\section[(dialogue)]{The house and describing words — a second pass}

\label{lesson:RU-R125-second-pass-the-house-and-describing-words}



\begin{quote}
Say the words for a shark and a nail.
\end{quote}

\subsection*{Your turn}
Say these aloud:

\begin{itemize}
\raggedright
  \item a shark, a nail, skin, a stomach, a throat
  \item a chin, a rainbow, lightning, thunder, a planet
  \item space, a bottle, a mirror, stairs, pink
  \item purple, brown, grey, orange, busy
  \item free, open, closed, quiet, noisy
\end{itemize}

\begin{tcolorbox}[breakable,colback=teal!4,colframe=teal!35!black,title={Before you move on}]
A shark? (\textbf{\ru{акула}}.) A nail? (\textbf{\ru{ноготь}}.) Skin? (\textbf{\ru{кожа}}.) A stomach? (\textbf{\ru{желудок}}.) A throat? (\textbf{\ru{горло}}.) A chin? (\textbf{\ru{подбородок}}.) A rainbow? (\textbf{\ru{радуга}}.) Lightning? (\textbf{\ru{молния}}.) Thunder? (\textbf{\ru{гром}}.) A planet? (\textbf{\ru{планета}}.) Space? (\textbf{\ru{космос}}.) A bottle? (\textbf{\ru{бутылка}}.) A mirror? (\textbf{\ru{зеркало}}.) Stairs? (\textbf{\ru{лестница}}.) Pink? (\textbf{\ru{розовый}}.) Purple? (\textbf{\ru{фиолетовый}}.) Brown? (\textbf{\ru{коричневый}}.) Grey? (\textbf{\ru{серый}}.) Orange? (\textbf{\ru{оранжевый}}.) Busy? (\textbf{\ru{занятый}}.) Free? (\textbf{\ru{свободный}}.) Open? (\textbf{\ru{открытый}}.) Closed? (\textbf{\ru{закрытый}}.) Quiet? (\textbf{\ru{тихий}}.) Noisy? (\textbf{\ru{шумный}}.)
\end{tcolorbox}
